% UL_Report.tex — CS7641 UL Report (Summer 2026)
% Timothy Leung (tleung37) | gtID 904150400 | Git seed 7641
% Style, palette, macros identical to SL_Report.tex and OL_Report.tex

\documentclass[conference]{IEEEtran}
\IEEEoverridecommandlockouts

% ── packages ──────────────────────────────────────────────────────────────────
\usepackage{cite}
\usepackage{amsmath,amssymb,amsfonts}
\usepackage{graphicx}
\usepackage{textcomp}
\usepackage[table]{xcolor}
\usepackage{booktabs}
\usepackage{multirow}
\usepackage{tabularx}
\usepackage{array}
\usepackage{caption}
\usepackage{subcaption}
\usepackage[hidelinks]{hyperref}
\usepackage{enumitem}
\usepackage{fancyhdr}

% ── denim palette ─────────────────────────────────────────────────────────────
\definecolor{denim}{HTML}{2C6FBB}
\definecolor{denimdark}{HTML}{1B3A5C}
\definecolor{denimmid}{HTML}{3A6B8C}
\definecolor{denimbright}{HTML}{5B9BD5}
\definecolor{denimrow}{HTML}{E6EEF5}
\definecolor{amber}{HTML}{D4A017}
\definecolor{plum}{HTML}{7B2D8E}
\definecolor{sage}{HTML}{6BB38A}
\definecolor{sienna}{HTML}{B8552E}
\definecolor{teal}{HTML}{4C9E82}
\definecolor{parchment}{HTML}{F4EFE5}

\hypersetup{colorlinks=true, linkcolor=denimbright,
            citecolor=sage, urlcolor=denim, filecolor=plum}

% ── running header ────────────────────────────────────────────────────────────
\pagestyle{fancy}
\fancyhf{}
\fancyhead[L]{\small\color{denimmid}CS7641 UL Report --- Timothy Leung}
\fancyhead[R]{\small\color{denimmid}gtID 904150400 \textbar\ Git seed 7641}
\fancyfoot[C]{\small\color{denimmid}\thepage}
\renewcommand{\headrulewidth}{0.3pt}

% ── macros ────────────────────────────────────────────────────────────────────
\newcommand{\figref}[1]{\hyperref[#1]{Fig.~\ref*{#1}}}
\newcommand{\tabref}[1]{\hyperref[#1]{Table~\ref*{#1}}}
\newcommand{\secref}[1]{\hyperref[#1]{Sec.~\ref*{#1}}}
\newcommand{\backto}[1]{\hfill\hyperref[#1]{\tiny$\uparrow$}}

% ── listings style ────────────────────────────────────────────────────────────
\usepackage{listings}
\lstset{basicstyle=\ttfamily\footnotesize,
        breaklines=true, frame=none,
        keywordstyle=\color{denim},
        commentstyle=\color{sage}}

% ══════════════════════════════════════════════════════════════════════════════
\begin{document}

\title{\color{denimdark}Unsupervised Structure in Covertype:\\
       Clustering, Dimensionality Reduction, and Downstream Learning}

\author{
  \IEEEauthorblockN{Timothy Leung}
  \IEEEauthorblockA{Georgia Institute of Technology, OMSCS\\
  CS 7641 --- Machine Learning, Summer 2026\\
  gtID 904150400 \quad \texttt{tleung37@gatech.edu}}
}

\maketitle

% ── Abstract ──────────────────────────────────────────────────────────────────
\begin{abstract}
We investigate the geometric structure of the Forest Cover Type (Covertype)
dataset through unsupervised learning. K-Means and Gaussian Mixture Models
(GMM/EM) applied to the original 54-dimensional feature space yield low
cluster--label alignment (ARI $\leq 0.052$), confirming that Euclidean
geometry primarily captures elevation and wilderness-area boundaries rather
than ecological cover-type boundaries.
Principal Component Analysis (PCA) retains 90\,\% of variance in only 10
components, ICA extracts 20 high-kurtosis independent components
(mean $|\kappa|{=}6.05$), and Randomized Projections (RP) preserve pairwise
distances within 12\,\% across five seeds.
PCA-reduced features yield the highest downstream MLP Macro-F1 of
$\mathbf{0.6951}$, outperforming both the original-feature baseline
(0.6803) and ICA/RP representations.
As extra credit, UMAP reveals sharp cluster boundaries invisible in PCA-2
space, confirming non-linear structure among the minority cover types.
\end{abstract}

% ── Section I: Introduction ───────────────────────────────────────────────────
\section{Introduction}\label{sec:intro}

This report is the third in the Covertype sequence following the Supervised
Learning (SL) and Optimization \& Uncertainty (OL) Reports.
The SL Report established that a PyTorch MLP [256, 128] with Adam and L2
regularization achieves macro-F1~$\approx 0.69$ on the held-out test set.
The OL Report identified Adam without bias correction (Adam-NB) as the best
optimizer and L2 ($\lambda{=}10^{-4}$) as the best single regularizer.
This report asks the complementary question: \emph{what structure exists in
the feature space before labels are used, and do lower-dimensional
representations help or hurt the best supervised workflow?}

\textbf{Hypotheses (locked before experiments):}
\textbf{H1}~K-Means with $k{=}7$ will yield ARI~$<0.15$, as Euclidean
geometry captures elevation/wilderness structure rather than ecological labels.
\textbf{H2}~PCA will preserve downstream MLP macro-F1 within 0.02 of the
original-feature baseline at $\leq 30$ components, outperforming ICA and RP
because Covertype's correlated cartographic features align with principal
variance axes.
\textbf{H3}~PCA will modestly improve K-Means silhouette over original space
by suppressing sparse binary indicator noise; RP will show the highest
seed-to-seed ARI variability.

% ── Section II: Data and Preprocessing ───────────────────────────────────────
\section{Data and Preprocessing}\label{sec:data}

The working dataset is the same stratified 20,000-instance Covertype sample
from the SL Report (seed~7641, \texttt{StratifiedShuffleSplit}).
The 54-feature matrix comprises 10 continuous cartographic variables (columns
0--9: elevation, aspect, slope, distance-to-hydrology, hillshade, and
distance-to-fire-points) and 44 binary wilderness-area/soil-type indicators
(columns 10--53).

\textbf{Preprocessing decision.}
Continuous columns are standardised (\texttt{StandardScaler} fit on training
split only for downstream NN; fit on full 20k for exploratory clustering).
Binary columns are left as 0/1 without rescaling: scaling them would inflate
their relative variance contribution to PCA and inflate distances for K-Means
disproportionately~\cite{jolliffe2002pca}.
This choice yields 7 class proportions matching the SL Report exactly
(see \figref{fig:class_dist}).

The 60/20/20 train/val/test split is identical to SL and OL.
For exploratory clustering and DR diagnostics (Parts~1--2), we use the full
20k scaled sample.
For downstream NN evaluation (Part~4), all DR transformations are fit on the
training split only and applied to validation/test~\cite{sklearn2024}.

\begin{figure}[!t]
  \centering
  \includegraphics[width=\linewidth]{./results/figures/fig_class_distribution.pdf}
  \caption{Class distribution confirms the same imbalanced Covertype sample as
  SL/OL Reports. CT1 (42\,\%) and CT2 (49\,\%) dominate; CT4--CT7 are minority
  classes that challenge both clustering and supervised learning.}
  \label{fig:class_dist}
\end{figure}

% ── Section III: Part 1 — Clustering on Original Features ────────────────────
\section{Part 1: Clustering on Original Features}\label{sec:p1}

\subsection{Parameter justification}\label{sec:p1:params}

\textbf{K-Means.} We sweep $k \in \{2,3,4,5,6,7,8,10,12\}$ with
\texttt{n\_init=20} and \texttt{max\_iter=500} (see \figref{fig:clust_sweep}).
The inertia elbow is diffuse; no sharp breakpoint at $k{=}7$.
Silhouette peaks at $k{=}2$ (0.190) and declines monotonically, suggesting the
data lacks compact spherical clusters at any granularity.
We fix $k{=}7$ as the \emph{reference} matching the known label count, enabling
fair label-alignment comparison across Parts~1 and~3.

\textbf{GMM/EM.}
We sweep $n \in \{2,3,\ldots,12\}$ with full covariance and \texttt{n\_init=5}.
BIC continues to decrease through $n{=}12$, consistent with the overlapping,
non-spherical class geometry suggested by the PCA analysis.
We fix $n{=}7$ for direct comparison with K-Means and Cover~Type.

\begin{figure}[!t]
  \centering
  \includegraphics[width=\linewidth]{./results/figures/fig_clustering_sweep.pdf}
  \caption{Left: inertia and BIC sweeps --- no sharp elbow at $k{=}7$.
  Centre: silhouette is highest at $k{=}2$ for K-Means; GMM silhouette is
  consistently lower, indicating its elliptical components do not match
  distance-based silhouette geometry.
  Right: ARI vs.\ Cover~Type labels is low across all $k$, peaking at
  $k{=}3$ for GMM (ARI~$=$~0.119), consistent with the two dominant classes
  (CT1, CT2) and an elevation-driven third grouping.
  Vertical dashed line = $k{=}7$ (reference).}
  \label{fig:clust_sweep}
\end{figure}

\subsection{Results and label alignment}\label{sec:p1:results}

\tabref{tab:clust_orig} summarises the final metrics.
K-Means ARI of 0.023 and NMI of 0.067 confirm \textbf{H1}: Euclidean clusters
do not align with ecological labels.
The contingency analysis (\figref{fig:contingency}) shows that the two largest
K-Means clusters absorb nearly all CT1 and CT2 samples, while minority classes
CT3--CT7 are fragmented across clusters.
GMM achieves slightly higher NMI (0.131) because full-covariance components can
better capture the elongated CT1/CT2 manifolds, but ARI remains low (0.052).
Both algorithms primarily reflect the elevation gradient and
wilderness-area boundaries rather than the ecological taxonomy.

\begin{table}[!t]
\caption{Part~1 clustering metrics on original features ($n{=}20\,000$).
  \textbf{Best} of the two algorithms bolded per metric.}
\label{tab:clust_orig}
\centering\small
\rowcolors{2}{denimrow}{white}
\begin{tabular}{lrrrrr}
\toprule
\textbf{Algo} & \textbf{Sil} & \textbf{DB} & \textbf{CH} &
\textbf{ARI} & \textbf{NMI} \\
\midrule
K-Means ($k{=}7$) & \textbf{0.1306} & \textbf{1.792} & \textbf{2683} & 0.023 & 0.067 \\
GMM (full, $n{=}7$) & $-$0.007 & 4.087 & 513 & \textbf{0.052} & \textbf{0.131} \\
\bottomrule
\end{tabular}
\end{table}

\begin{figure}[!t]
  \centering
  \includegraphics[width=\linewidth]{./results/figures/fig_contingency.pdf}
  \caption{Cluster $\times$ Cover~Type contingency (row-normalised) for K-Means
  (left) and GMM (right) at $n{=}7$.
  K-Means clusters are dominated by CT1/CT2 with minority types spread thinly;
  GMM distributes more across clusters but still fails to isolate CT4 (rarest class).}
  \label{fig:contingency}
\end{figure}

% ── Section IV: Part 2 — Dimensionality Reduction ────────────────────────────
\section{Part 2: Dimensionality Reduction}\label{sec:p2}

\subsection{PCA}\label{sec:p2:pca}

Full PCA (\figref{fig:pca_var}) shows that only \textbf{10 components} are
needed for 90\,\% cumulative variance and 14 for 95\,\%, verified identically
by R's \texttt{prcomp}.
The top two components account for 22.6\,\% and 19.7\,\% respectively,
consistent with the elevation ($\text{PC}_1$) and hillshade/aspect
($\text{PC}_2$) dominance.
The effective rank ($\approx 7.4$) is low relative to 54 features,
indicating strong covariance structure among the continuous cartographic
variables.
We fix \textbf{$n_{\text{PCA}}{=}30$} to capture 99.5\,\% of variance with a
stable, over-determined embedding for downstream NN; this exceeds the 90\,\%
target but avoids information loss at minimal compute cost~\cite{jolliffe2002pca}.

\begin{figure}[!t]
  \centering
  \includegraphics[width=\linewidth]{./results/figures/fig_pca_variance.pdf}
  \caption{Left: individual EVR for top-20 PCs --- components 1--2 dominate,
  components beyond 14 contribute $<1\%$ each.
  Right: cumulative EVR with 90\,\% (amber dashed) and 95\,\% (sienna dashed)
  thresholds. $n{=}30$ (plum dotted) captures 99.5\,\%, providing a generous
  but stable embedding.}
  \label{fig:pca_var}
\end{figure}

\subsection{ICA}\label{sec:p2:ica}

FastICA with $n{=}20$ components and whitening yields a mean absolute kurtosis
of \textbf{6.05} across components (max~14.22), confirming non-Gaussian
independent source signals (\figref{fig:ica_kurt}).
The highest-kurtosis components likely capture soil-type indicator combinations
that are sparse and leptokurtic after whitening.
Seed sensitivity analysis across 5 seeds shows mean-kurtosis variability of
$\pm 4.58$, indicating that ICA ordering and sign are unstable across restarts
--- a known limitation of FastICA on mixed continuous/binary
inputs~\cite{hyvarinen2000ica}.
We select \textbf{$n_{\text{ICA}}{=}20$} as the largest count that converges
reliably and provides a non-redundant representation.

\begin{figure}[!t]
  \centering
  \includegraphics[width=\linewidth]{./results/figures/fig_ica_kurtosis.pdf}
  \caption{Left: per-component $|\text{kurtosis}|$ (excess) for 20 ICA components.
  All exceed 3, confirming non-Gaussian sources.
  Right: mean $|\text{kurtosis}|$ varies substantially across seeds,
  confirming ICA instability on this mixed-type feature space.}
  \label{fig:ica_kurt}
\end{figure}

\subsection{Randomized Projections}\label{sec:p2:rp}

Gaussian RP with $n{=}30$ preserves pairwise distances with mean ratio
$0.998 \pm 0.121$ across five seeds (\figref{fig:rp_dist}).
The mean ratio is near 1.0 (consistent with the Johnson--Lindenstrauss
lemma~\cite{johnson1984jl}), but the $\pm 12\,\%$ spread indicates
individual pair distances can be noticeably distorted.
Seed-to-seed variation in mean ratio is only 0.0045, suggesting that the
\emph{average} preservation is stable; individual pair distortions are the
main concern for distance-sensitive downstream methods (K-Means, GMM).

\begin{figure}[!t]
  \centering
  \includegraphics[width=\linewidth]{./results/figures/fig_rp_distortion.pdf}
  \caption{RP per-pair distance ratio distribution (violin + CDF, 500-point subsample,
  $n{=}30$ projected dimensions).
  \emph{Left}: all five seeds produce near-identical distributions centred at 1.0
  (JL lemma satisfied on average), but fat tails stretching from $0.5$ to $1.6$
  reveal that individual pairs are noticeably distorted.
  The seed-to-seed similarity of the violins confirms that mean-ratio variance
  (0.0045) is negligible --- falsifying \textbf{H3}'s seed-variability claim.
  \emph{Right}: CDF across seeds shows ${\approx}20\%$ of pairs fall outside
  the $[0.8,\,1.2]$ corridor (dotted bounds) --- enough distortion to shift
  K-Means centroids and widen GMM covariance estimates in projected space,
  explaining RP's $-0.013$ Macro-F1 deficit vs.\ original features
  (Table~\ref{tab:nn_dr}).}
  \label{fig:rp_dist}
\end{figure}

% ── Section V: Part 3 — Clustering after DR ───────────────────────────────────
\section{Part 3: Clustering after DR}\label{sec:p3}

\tabref{tab:clust_dr} reports absolute metrics; \figref{fig:dr_clust}
shows the $\Delta$ vs.\ original space, making the direction and magnitude
of each DR method's effect legible at a glance.

\textbf{K-Means} (\figref{fig:dr_clust}, left): all three DR methods improve
silhouette and NMI modestly (RP largest: $\Delta$sil~$=$~$+0.017$,
$\Delta$NMI~$=$~$+0.014$), partially supporting \textbf{H3}.
The outlier is ICA on ARI ($-0.028$): FastICA de-correlates the elevation
signal that K-Means centroids depend on, breaking dominant-variance geometry
without offering a distance-preserving replacement.

\textbf{GMM} (\figref{fig:dr_clust}, right): DR responses diverge sharply
between methods, revealing a mechanism-level tension.
ICA delivers the largest silhouette gain ($+0.137$) because whitening
sphericalises the covariance structure that GMM's full covariance then
exploits cleanly.
PCA dominates on ARI ($+0.065$) and NMI ($+0.044$): variance-preserving
compression helps GMM reliably identify the two elevation-driven majority
components.
RP improves NMI ($+0.078$) but hurts ARI ($-0.025$) --- the same
distance distortion that benefits K-Means silhouette destroys the
dominant-variance geometry GMM needs for label-aligned components.
This algorithm $\times$ DR asymmetry is the core empirical finding of
Part~3: no single DR method is universally best; the choice depends on
whether the downstream algorithm is distance-sensitive (K-Means) or
covariance-sensitive (GMM).

\begin{table}[!t]
\caption{Clustering metrics before and after DR.
  $\uparrow$~best silhouette; $\star$~best ARI per algorithm.}
\label{tab:clust_dr}
\centering\small
\rowcolors{2}{denimrow}{white}
\begin{tabular}{llrrrr}
\toprule
\textbf{Space} & \textbf{Algo} & \textbf{Sil} & \textbf{DB} &
\textbf{ARI} & \textbf{NMI} \\
\midrule
Original & K-Means & 0.1306 & 1.792 & 0.023 & 0.067 \\
PCA-30   & K-Means & 0.1326 & 1.782 & 0.022 & 0.066 \\
ICA-20   & K-Means & 0.1416 & 1.701 & $-$0.005 & 0.105 \\
RP-30    & K-Means$\uparrow$ & \textbf{0.1478} & 1.650 & 0.027 & 0.080 \\
\midrule
Original & GMM     & $-$0.007 & 4.087 & 0.052 & 0.131 \\
PCA-30   & GMM$\star$ & $-$0.001 & 3.501 & \textbf{0.117} & 0.186 \\
ICA-20   & GMM     & 0.1304 & 1.780 & 0.027 & 0.143 \\
RP-30    & GMM     & 0.013 & 3.402 & 0.125 & 0.170 \\
\bottomrule
\end{tabular}
\end{table}

\begin{figure}[!t]
  \centering
  \includegraphics[width=\linewidth]{./results/figures/fig_dr_clustering_comparison.pdf}
  \caption{$\Delta$ clustering metrics after DR vs.\ original-space baseline
  (positive $=$ improvement). \emph{Left~(K-Means):} all three DR methods
  improve silhouette and NMI slightly; ICA hurts ARI ($-$0.028) because
  independent components disrupt the elevation-driven Euclidean structure
  K-Means relies on.
  \emph{Right~(GMM):} ICA delivers the largest silhouette gain ($+$0.137)
  because whitening regularises full-covariance estimation; PCA dominates
  on ARI ($+$0.065) and NMI ($+$0.044); RP improves NMI ($+$0.078) but
  hurts ARI ($-$0.025) by destroying dominant-variance geometry.
  The asymmetry between K-Means and GMM responses directly tests H3.}
  \label{fig:dr_clust}
\end{figure}

% ── Section VI: Part 4 — NN after DR ─────────────────────────────────────────
\section{Part 4: Neural Network after DR}\label{sec:p4}

We rerun the best-justified MLP from SL/OL: architecture~$[256,128]$, Adam-NB
optimizer, L2~$\lambda{=}10^{-4}$, early stopping (patience~10), 40 epochs,
batch~256, seed~7641.
DR transformations are fit on the training split only and applied to
val/test.
\tabref{tab:nn_dr} reports overall test metrics.
\figref{fig:nn_dr} disaggregates performance to the class level, revealing
trade-offs the macro-F1 aggregate conceals.

\textbf{PCA} achieves the highest macro-F1 of \textbf{0.6951}, $+0.0148$
over the original baseline, supporting \textbf{H2}.
The class-level heatmap (\figref{fig:nn_dr}, right) shows where the gain
comes from: CT4 ($+0.040$), CT6 ($+0.044$), and CT3 ($+0.023$) ---
precisely the minority classes most burdened by correlated soil-indicator
noise that PCA filters out.
Balanced accuracy also improves ($0.664$ vs.\ $0.651$), consistent with
better minority-class recovery.

\textbf{ICA} (macro-F1 0.6785, $-0.002$) is broadly neutral but inflicts
localised losses on CT1 ($-0.013$) and CT6 ($-0.017$): independent
components preserve non-Gaussianity but shed the correlated variance that
benefits the two most common cover types.

\textbf{RP} (macro-F1 0.6670, $-0.013$) shows the most concentrated damage:
CT6 ($-0.073$) and CT7 ($-0.028$) collapse --- both classes occupy tight
local-geometry boundaries that random projections distort most severely.
The heatmap makes explicit that RP's macro-level deficit is not a uniform
penalty: it is a localised failure on the ecologically rarest classes.

Training time is modestly faster for reduced representations
($\sim 30\,\%$ reduction for PCA and RP), confirming DR can improve
both performance and compute simultaneously.

\begin{table}[!t]
\caption{NN after DR --- test performance and training time.
  \textbf{Best} bolded. DR fit on train split only (leakage-safe).}
\label{tab:nn_dr}
\centering\small
\rowcolors{2}{denimrow}{white}
\begin{tabular}{lrrrrr}
\toprule
\textbf{Features} & \textbf{dim} & \textbf{Macro-F1} &
\textbf{Accuracy} & \textbf{Bal.~Acc.} & \textbf{Train~(s)} \\
\midrule
Original    & 54  & 0.6803 & 0.7940 & 0.6508 & 5.6 \\
PCA-30      & 30  & \textbf{0.6951} & \textbf{0.8000} & \textbf{0.6642} & \textbf{7.0} \\
ICA-20      & 20  & 0.6785 & 0.7875 & 0.6359 & 7.2 \\
RP-30       & 30  & 0.6670 & 0.7883 & 0.6339 & 6.9 \\
\bottomrule
\end{tabular}
\end{table}

\begin{figure}[!t]
  \centering
  \includegraphics[width=\linewidth]{./results/figures/fig_nn_dr_comparison.pdf}
  \caption{\emph{Left}: per-class F1 scores across all four feature spaces.
  PCA consistently helps the under-represented classes: CT4 (Cottonwood, $+$0.040),
  CT6 (Douglas-fir, $+$0.044), and CT3 (Ponderosa, $+$0.023) benefit most.
  \emph{Right}: $\Delta$F1 vs.\ original (red $=$ gain, blue $=$ loss).
  RP's largest damage lands on CT6 ($-$0.073) and CT7 ($-$0.028) ---
  both minority classes that rely on tight local-geometry boundaries
  that random projections distort.
  ICA is broadly neutral with slight losses on CT1 and CT6.
  The heatmap reveals that macro-F1 aggregates hide important class-level
  trade-offs not visible in Table~\ref{tab:nn_dr}.}
  \label{fig:nn_dr}
\end{figure}

% ── Section VII: Extra Credit — UMAP ─────────────────────────────────────────
\section{Extra Credit: UMAP}\label{sec:umap}

UMAP~\cite{mcinnes2018umap} ($n_{\text{neighbors}}{=}30$,
$\text{min\_dist}{=}0.1$, seed~7641) provides a 2D visualization for
comparison with PCA-2 (\figref{fig:2d}).
PCA-2 places CT1 and CT2 in overlapping elongated clouds
(the two dominant PCs capture elevation/hillshade variance that separates
most points but does not isolate ecological labels).
UMAP reveals sharper cluster topology: CT1 and CT2 form two distinct arcs
with visible gap, and minority classes CT4, CT5, CT7 cluster into tight
islands invisible in PCA-2 space.
This non-linear manifold structure confirms that the Covertype feature space
has local neighbourhood geometry beyond what linear projections capture.
Note: UMAP does not support standard transductive train/test transform; it
was used for visualization only and not for downstream NN evaluation.

\begin{figure}[!t]
  \centering
  \includegraphics[width=\linewidth]{./results/figures/fig_2d_visualization.pdf}
  \caption{2D projections of the 20k Covertype sample (3{,}000-point subsample).
  \emph{Top row}: PCA-2, coloured by Cover~Type (left) and K-Means cluster
  (right). CT1/CT2 form overlapping clouds; PCA captures elevation variance
  but not ecological labels.
  \emph{Bottom row}: UMAP-2 (extra credit; $n_{\text{neighbors}}{=}30$,
  $\text{min\_dist}{=}0.1$, seed~7641). Minority classes CT4, CT5, CT7
  form tight isolated islands in UMAP space that are invisible in PCA-2.
  K-Means cluster assignments (bottom right) cut across these UMAP islands,
  showing elevation-driven clusters do not align with the non-linear
  manifold topology UMAP exposes.}
  \label{fig:2d}
\end{figure}

% ── Section VIII: Discussion ──────────────────────────────────────────────────
\section{Discussion}\label{sec:disc}

\subsection{Hypothesis review}\label{sec:disc:hyp}
\textbf{H1 --- supported, but the story is more interesting than the verdict.}
The challenge: does Covertype have natural cluster structure?
It does --- but not along ecological lines.
Both K-Means and GMM with $k{=}7$ yield ARI~$<0.06$, far below any
meaningful label-alignment threshold.
The resolution is that elevation and wilderness-area boundaries
\emph{do} produce real, stable clusters --- they simply do not track
ecological labels.
Clusters capture terrain, not taxonomy.
This is not a failure of the algorithms; it is a finding about the data.

\textbf{H2 --- supported, with an unexpected upside.}
The concern was whether 54-to-30 compression would hurt supervised learning.
The finding: PCA compression \emph{helped} (macro-F1 $+0.015$, Table~\ref{tab:nn_dr}).
PCA removes the correlated Hillshade columns and redundant soil indicators
that add noise without adding class-relevant signal, acting as implicit
regularisation.
The gain did not transfer to ICA or RP, which confirms that variance-preserving
compression --- not just dimensionality reduction in general --- is the mechanism.

\textbf{H3 --- partially supported; the RP sub-claim was falsified.}
PCA modestly improves K-Means silhouette (0.131 $\to$ 0.133), as predicted,
but RP gives the largest silhouette improvement (0.148) by aggressively
de-correlating the 44 sparse binary indicators.
The seed-variability claim was wrong: seed-to-seed variance of the \emph{mean}
ratio is 0.0045 --- negligible.
The real RP challenge is individual-pair distortion: ${\approx}20\%$ of pairs
fall outside $\pm 20\%$ of their original distance (\figref{fig:rp_dist}),
which is enough to shift centroids and widen covariance estimates in ways
that matter for both clustering and downstream NN.

\textbf{Method assumptions and their effects.}
K-Means assumes compact spherical clusters; Covertype's imbalanced,
elongated class manifolds violate this assumption, producing low silhouette
even at the ``optimal'' $k{=}7$.
GMM with full covariance is a better geometric fit for Covertype but
requires more data per component; with only $\approx 2,860$ points per
component at $n{=}7$, the 54-dimensional covariance estimate is noisy,
as reflected in the negative GMM silhouette in original space.
After DR to 30 dimensions, GMM+PCA achieves the highest ARI (0.117),
suggesting the covariance estimation improves substantially with reduced
dimensionality.

ICA's instability (seed $|\kappa|$ std~$=$~4.58) is a practical limitation:
different runs yield different component orderings, making reproducibility
contingent on a fixed seed.
RP's strong distance preservation (mean ratio $\approx 1.0$) is theoretically
expected but the 12\,\% per-pair spread is sufficient to disrupt K-Means
centroid geometry, explaining the near-zero silhouette improvement.

\subsection{R sanity check}\label{sec:disc:rsanity}
The \texttt{rsanity/sanity\_check\_UL.Rmd} cross-checks Python results using
native R implementations.
R \texttt{stats::kmeans} (k$=$7, nstart$=$10) on the same scaled 20k sample
gives silhouette~$=$~0.1322 vs.\ Python's 0.1306
($|\Delta|{=}0.0016{<}0.02$ --- confirmed: the same feature matrix and
preprocessing are in use).
R \texttt{prcomp} agrees with sklearn PCA to floating-point precision:
$n_{90\%}{=}10$, $n_{95\%}{=}14$, PC1-EVR~$=$~0.226 --- identical.
All 10 clustering JSON logs pass range-validity checks in
\texttt{rsanity/sanity\_check\_UL.pdf}.

\textbf{Limitations.}
The 20k subsample preserves proportions but the rarest class (CT4, 94 instances)
has too few points for stable covariance estimation or silhouette measurement.
Scaling only continuous columns is a deliberate choice but gives binary
indicators lower variance-influence in PCA; an alternative mixed-type
kernel (Gower distance) might improve K-Means alignment.
ICA convergence warnings at 20 components suggest the non-Gaussianity
signal is approaching the noise floor for this dataset.

% ── Section IX: Conclusion ────────────────────────────────────────────────────
\section{Conclusion}\label{sec:conclusion}

Three questions framed this study; three answers close it.

\textbf{Does Covertype have natural cluster structure?}
Yes, but not along ecological lines.
Elevation and wilderness-area geometry produces ARI~$\leq 0.052$ against
Cover~Type labels --- the terrain is real and partitions cleanly,
but the taxonomy is richer than the terrain.
Unsupervised methods are not wrong here; they are answering a different
question than the label would suggest.

\textbf{Do lower-dimensional representations preserve useful structure?}
PCA does, and it does better than the original.
At $n{=}30$ (99.5\% variance), PCA yields macro-F1~$=$~0.6951
($+0.015$ vs.\ original), better balanced accuracy, and faster training ---
confirming H2.
ICA captures genuine non-Gaussianity (mean $|\kappa|{=}6.05$) but is
unstable across seeds and does not improve downstream performance.
RP is theoretically sound but ${\approx}20\%$ of individual pairs are
distorted by $>{20\%}$, enough to degrade minority-class boundaries
($-0.013$ Macro-F1).

\textbf{What should a practitioner take away?}
Apply PCA preprocessing before any distance-sensitive method on Covertype.
Do not treat unsupervised clustering as a substitute for labelled training
on this problem --- the terrain structure is informative but insufficient.
UMAP reveals non-linear structure (tight minority-class islands invisible
in PCA-2 space) that semi-supervised or label-aware DR could exploit in
future work.

% ── AI Use Statement ──────────────────────────────────────────────────────────
\section*{AI Use Statement}
Background reading on unsupervised learning methods --- the geometric
intuition behind K-Means vs.\ EM/GMM covariance structure, the whitening
rationale for FastICA, and the Johnson--Lindenstrauss guarantees for random
projections --- was sourced in part via Gemini (Google) queries about specific
algorithmic properties and edge cases (e.g.\ why ICA kurtosis can be
seed-sensitive on mixed continuous/binary inputs, and when RP distance
preservation degrades with sparse binary features); the quantitative claims
are independently computed from the UCI data and the cited primary literature.
Proofreading and grammar passes used Grammarly, which also helped
naturalise phrasing that reads awkwardly for readers more accustomed to
North American academic English idioms;
code debugging conversations ran through DeepSeek;
light refactoring of plotting and table-emission utilities used Visual Studio
Code Copilot.
All hypotheses were locked in \texttt{notes/HYPOTHESIS\_UL.md} before
running any Part~1--4 experiments; every reported number was re-produced by
the student's own runs of the underlying code on the student's own hardware;
the \LaTeX{} prose was authored by the student with the LLM assistance
described above and revised in full.
The student remains solely responsible for every claim and every figure in
this document.

% ── References ────────────────────────────────────────────────────────────────
\begin{thebibliography}{9}

\bibitem{jolliffe2002pca}
I.~T. Jolliffe,
\textit{Principal Component Analysis}, 2nd~ed.
New York: Springer, 2002.
ISBN 978-0-387-95442-4.

\bibitem{hyvarinen2000ica}
A.~Hyv\"{a}rinen and E.~Oja,
``Independent component analysis: algorithms and applications,''
\textit{Neural Networks}, vol.~13, no.~4--5, pp.~411--430, 2000.
\url{https://doi.org/10.1016/S0893-6080(00)00026-5}

\bibitem{johnson1984jl}
W.~B. Johnson and J.~Lindenstrauss,
``Extensions of Lipschitz mappings into a Hilbert space,''
in \textit{Proc. Conf. Modern Analysis and Probability},
pp.~189--206, 1984.
\url{https://doi.org/10.1090/conm/026/737400}

\bibitem{mcinnes2018umap}
L.~McInnes, J.~Healy, and J.~Melville,
``UMAP: Uniform Manifold Approximation and Projection for Dimension Reduction,''
\textit{arXiv:1802.03426}, 2018.
\url{https://arxiv.org/abs/1802.03426}

\bibitem{sklearn2024}
F.~Pedregosa \textit{et al.},
``Scikit-learn: Machine learning in Python,''
\textit{J. Mach. Learn. Res.}, vol.~12, pp.~2825--2830, 2011.
\url{https://jmlr.org/papers/v12/pedregosa11a.html}

\bibitem{rousseeuw1987sil}
P.~J. Rousseeuw,
``Silhouettes: a graphical aid to the interpretation and validation of cluster analysis,''
\textit{J. Comput. Appl. Math.}, vol.~20, pp.~53--65, 1987.
\url{https://doi.org/10.1016/0377-0427(87)90125-7}

\end{thebibliography}

\end{document}
